\documentclass[
    language=english,
    doctype=article,
    institution=none,
    titlestyle=book,
]{../../omnilatex}

\RequirePackage{config/document-settings}
\addbibresource{../../bib/bibliography.bib}

\begin{document}

\maketitle

\section*{Abstract}
Natural language processing (NLP) has advanced rapidly for high-resource
languages, yet over 7000 languages worldwide remain under-served by modern
computational tools.  This literature review systematically examines transfer
learning techniques that adapt models trained on resource-rich languages to
low-resource settings.  We organise the reviewed literature into four thematic
strands: cross-lingual pre-training, few-shot and zero-shot adaptation,
data augmentation strategies, and evaluation frameworks.  The review identifies
critical gaps in evaluation methodology and highlights promising directions for
future research.

\section{Introduction}
The success of large pretrained language models such as BERT, GPT, and
XLM-R has transformed NLP for English and other high-resource
languages~\cite{bibid}.  However, the vast majority of the world's languages
lack sufficient annotated data to train comparable models from
scratch~\cite{einstein1905}.  Transfer learning offers a principled approach to
bridging this resource gap by leveraging knowledge encoded in models trained on
data-rich languages and adapting it to target languages with limited resources.

This review addresses the following research questions:
\begin{enumerate}
    \item What transfer learning methods are most effective for low-resource
          NLP tasks?
    \item How do cross-lingual representations generalise across typologically
          diverse languages?
    \item What evaluation standards are appropriate for assessing model
          performance in low-resource settings?
\end{enumerate}

\section{Cross-Lingual Pre-Training}

\subsection{Multilingual Language Models}
The introduction of multilingual BERT (mBERT) demonstrated that a single model
can learn representations transferable across 104 languages.  Subsequent work
extended this approach to over 100 languages with XLM-RoBERTa, achieving
state-of-the-art results on cross-lingual benchmarks~\cite{goossensLaTeXCompanion1993}.

\subsection{Language-Adaptive Pre-Training}
A complementary strategy involves continued pre-training on monolingual corpora
in the target language.  This domain-adaptive approach has shown improvements
of 3--7 F1 points on named entity recognition tasks for languages such as
Igbo, Yoruba, and Hausa.

\section{Few-Shot and Zero-Shot Adaptation}

\subsection{Prompt-Based Methods}
Large language models can be prompted to perform tasks without explicit
fine-tuning.  Techniques such as pattern-exploiting training (PET) and
contextual prompting achieve competitive results with as few as 32 labelled
examples per class.

\subsection{Cross-Lingual Transfer}
Zero-shot cross-lingual transfer involves fine-tuning a multilingual model on
a high-resource source language and directly applying it to a target language.
Performance correlates strongly with typological similarity between the source
and target language pair.

\section{Data Augmentation Strategies}

\subsection{Translation-Based Augmentation}
Translating labelled data from a high-resource language into the target
language is a straightforward augmentation strategy.  However, translation
artefacts and domain mismatch can introduce noise that degrades model
performance.

\subsection{Back-Translation and Paraphrase Generation}
Generating synthetic training examples through back-translation or neural
paraphrase models has proven effective, particularly for morphologically rich
languages where inflectional variation is substantial.

\section{Evaluation Frameworks}

\subsection{Benchmark Datasets}
Several benchmarks have been developed for evaluating low-resource NLP,
including AmericasNLI, AfriQA, and XTREME.  These datasets cover diverse
language families and task types.

\subsection{Evaluation Challenges}
A significant challenge in evaluating low-resource models is the scarcity of
held-out test sets.  Cross-validation and stratified sampling strategies are
commonly employed but may overestimate true out-of-sample performance.

\section{Synthesis and Gaps}
The reviewed literature reveals several persistent gaps:
\begin{description}
    \item[Script diversity] Most cross-lingual models assume shared scripts or
    transliteration into Latin characters, limiting applicability for languages
    using non-Latin scripts.
    \item[Oral languages] Languages with limited written traditions are largely
    absent from the transfer learning literature.
    \item[Evaluation standards] There is no consensus on appropriate evaluation
    protocols for extremely low-resource settings (fewer than 1000 labelled
    sentences).
\end{description}

\section{Conclusion}
Transfer learning has made substantial progress in extending NLP capabilities to
low-resource languages.  Future research should prioritise inclusive evaluation
methodologies, oral language processing, and community-driven dataset
development to ensure equitable advances across the world's linguistic
diversity.

\printbibliography

\end{document}
